\subsubsection{局部化与正则化}\label{subsec:distribution-localization}

虽然不能计算广义函数在点处的值, 但可以将它们限制到开集上. 这相当直接. 若 $G$ 是 $\Omega$ 的开子集, 则 $\mathcal{D}(G)$ 自然嵌入 $\mathcal{D}(\Omega)$, 因而 $\Omega$ 上的每个广义函数通过限制都定义了 $G$ 上的广义函数. 因此作如下定义.

\begin{Definition}
\label{def:distribution-vanishing}
若对每个 $\phi\in\mathcal{D}(G)$ 都有 $(f,\phi)=0$, 则称 $f\in\mathcal{D}^{\prime}(\Omega)$ 在开集 $G\subset\Omega$ 上\textbf{为零}. 若两个分布之差在 $G$ 上为零, 则它们在 $G$ 上相等.
\end{Definition}

可以证明 (参见 Problem~\ref{ex:distribution-local-vanishing}), 若 $f$ 在 $G$ 中每一点附近局部为零, 即 $G$ 的每一点都有一个使 $f$ 为零的邻域, 则 $f$ 在 $G$ 上为零. 一个直接推论是, 若 $f$ 在一族开集中的每一个上都为零, 则它在这些开集的并上也为零. 因而存在最大的开集 $N_f$, 使 $f$ 在其上为零.

\begin{Definition}
\label{def:distribution-support}
$N_f$ 在 $\Omega$ 中的补集称为 $f$ 的\textbf{支集}.
\end{Definition}

\begin{Example}
\label{ex:delta-support}
delta 函数的支集为集合 $\{0\}$. 虽然不能计算 delta 函数在点处的值, 但说它除原点外均为零是有意义的.
\end{Example}

\begin{Remark}
\label{rem:singular-function-regularization}
具有不可积奇性的函数不能通过方程 \eqref{eq:regular-distribution-pairing} 定义为广义函数. 然而, 通常可以定义一个广义函数, 它在任意不包含奇点的开集上均与该奇异函数一致. 这样的广义函数称为正则化. 例如, $\mathbb{R}$ 上函数 $1/x$ 的一个正则化由主值积分给出
\begin{equation}\label{eq:principal-value-one-over-x-regularization}
\begin{aligned}
(f,\phi)&=\int_{-\infty}^{-\epsilon}\dfrac{\phi(x)}x\,\mathrm{d}x
+\int_{-\epsilon}^{\epsilon}\dfrac{\phi(x)-\phi(0)}x\,\mathrm{d}x\\
&\qquad+\int_\epsilon^\infty\dfrac{\phi(x)}x\,\mathrm{d}x
\end{aligned}
\end{equation}
(参见 Problem~\ref{ex:principal-value-regularization}).
\end{Remark}
